\chapter{Dependency graph}

This blueprint records the public theorem graph of the orthogonal-row
formalization. It is not a substitute for \texttt{SPEC.md}. Status line
after Milestone D: machine-checked randomized quasi-optimal column
selection for the orthogonal-row case of Simons Problem 4.1. Do not say
Problem 4.1 is solved.

\section{Milestones A--D (closed)}

\begin{theorem}[Orthogonal-row / column correspondence]
\label{thm:milestone-a}
\lean{StructuredColumnSelection.milestoneA_transpose_correspondence}
\leanok
Working with $A\in\RR^{k\times n}$ and $AA^\top=I$ is equivalent to
working with $Q=A^\top$ and $Q^\top Q=I$. Selected squares and Grams are
available as definitions.
\end{theorem}

\begin{theorem}[Volume normalisation]
\label{thm:milestone-b}
\lean{StructuredColumnSelection.milestoneB_volume_normalization, StructuredColumnSelection.milestoneB_cauchyBinet}
\leanok
\uses{thm:milestone-a}
If $AA^\top=I$, then $\sum_{|J|=k}\det(A_J)^2=1$.
\end{theorem}

\begin{theorem}[Inverse-Gram identity]
\label{thm:milestone-c}
\lean{StructuredColumnSelection.milestoneC_inverse_gram_expectation, StructuredColumnSelection.milestoneC_adjugate_sum}
\leanok
\uses{thm:milestone-b}
$\sum_{|J|=k}\operatorname{adj}(A_J A_J^\top)=(n-k+1)I$ when $AA^\top=I$.
\end{theorem}

\begin{theorem}[Inverse-Frobenius expectation and Markov tail]
\label{thm:milestone-d}
\lean{StructuredColumnSelection.milestoneD_expected_inv_frob_sq, StructuredColumnSelection.milestoneD_markov_inv_frob}
\leanok
\uses{thm:milestone-c}
The volume-weighted inverse-Frobenius energy is $k(n-k+1)$, with a finite
Markov tail on that proxy.
\end{theorem}

\section{Milestone E, structural CPQR}

\begin{theorem}[CPQR full-rank stopping]
\label{thm:e-structural}
\lean{StructuredColumnSelection.milestoneE_leverage_sum, StructuredColumnSelection.milestoneE_first_pivot_is_max, StructuredColumnSelection.milestoneE_first_leverage_ge, StructuredColumnSelection.milestoneE_cpqr_card_eq}
\leanok
\uses{thm:milestone-a}
Empty residuals are leverages and sum to $k$. A first pivot maximises
empty residual, hence leverage $\ge k/n$. Orthogonal-row $\CPQR$ returns
exactly $k$ columns. This is not an inverse-norm bound.
\end{theorem}

\begin{theorem}[$k=1$ volume]
\label{thm:e-k1}
\lean{StructuredColumnSelection.milestoneE_k1_volume_ge}
\leanok
\uses{thm:e-structural}
For a single orthonormal row, $\CPQR$ volume is at least $n^{-1}$, matching
the workshop scale $\sqrt{n}$. This does not close Milestone E.
\end{theorem}

\begin{theorem}[$k=2$ inverse-trace]
\label{thm:e-k2}
\lean{StructuredColumnSelection.milestoneE_k2_inv_trace_le}
\leanok
\uses{thm:e-bidiagonal}
For two orthonormal rows, Gram--Schmidt $U$ is automatically bidiagonal, so
$\tr((G')^{-1})\le 3(n-1)$. This matches a joint polynomial bound for
$k=2$. It does not close Milestone E: SPEC §16 outcome 1 needs every $k$.
\end{theorem}

\begin{theorem}[Residual energy and binomial volume]
\label{thm:e-residual}
\lean{StructuredColumnSelection.milestoneE_residual_energy, StructuredColumnSelection.milestoneE_next_residual_ge, StructuredColumnSelection.milestoneE_cpqr_volume_ge_binomial}
\leanok
\uses{thm:e-structural}
Independent residuals sum to $k-\#J$. $\CPQR$ volume is at least
$1/\binom{n}{k}$. That factor is exponential in $k$ when $n\approx 2k$ and
is not a joint $\poly(n,k)$ bound.
\end{theorem}

\begin{theorem}[Triangular inverse-trace bound]
\label{thm:e-triangular}
\lean{StructuredColumnSelection.milestoneE_pivot_gram_inv_trace_le}
\leanok
\uses{thm:e-residual}
$\tr((G')^{-1})\le (n-k+1)(4^k+6k-1)/9$. Polynomial in $n$, exponential in
$k$. Not a joint $\poly(n,k)$ bound.
\end{theorem}

\begin{theorem}[Bidiagonal-$U$ inverse-trace]
\label{thm:e-bidiagonal}
\lean{StructuredColumnSelection.milestoneE_bidiagonal_U_inv_trace_le, StructuredColumnSelection.milestoneE_gsR_mul_transpose}
\leanok
\uses{thm:e-triangular}
If Gram--Schmidt $U$ is bidiagonal, then
$\tr((G')^{-1})\le k(k+1)/2\cdot(n-k+1)$. This is a hypothesis on $U$, not a
class of $A$, and not Path 1. Implicit $R$ satisfies $RR^\top=I$.
\end{theorem}

\begin{conjecture}[Milestone E remains open]
\label{conj:e-open}
\notready
\uses{thm:e-k1, thm:e-k2, thm:e-residual, thm:e-triangular, thm:e-bidiagonal}
A joint $\poly(n,k)$ inverse-norm or inverse-Frobenius bound for every
orthogonal-row $\CPQR$ set, or a certified superpolynomial counterexample,
or a counterexample plus a stronger static class (SPEC §16). Path 1 is not
closed for general $k$. Multi-neighbor coupling in $U$ still blocks the
general statement; Businger--Golub $|S_{ij}|\le 2^{j-i-1}$ is sharp for
arbitrary unit-triangular matrices.
\end{conjecture}

\section{Milestone F}

\begin{conjecture}[CSSP perturbation bridge]
\label{conj:f-open}
\notready
Do not start Milestone F until E is closed or the user asks. The intended
statement is $\|\Pi_J A-A\|\le(1+\|Q_J^{-1}\|)\|A-A_k\|$ for $A=B+E$ with
$B$ of rank $k$ and $Q$ the right singular frame of $B$.
\end{conjecture}
